\documentclass[11pt]{article}

\usepackage[T1]{fontenc}
\usepackage{lmodern}
\usepackage{amsmath,amssymb,amsthm}
\usepackage[a4paper,margin=1in]{geometry}
\usepackage[hidelinks]{hyperref}
\usepackage{microtype}

\newtheorem{proposition}{Proposition}
\theoremstyle{remark}
\newtheorem{remark}{Remark}

\title{Direct Generation of a Somos--4 Sequence\\
from an Algebraic Generating Function}
\author{Thomas Scheuerle}
\date{}

\begin{document}
\maketitle

\begin{abstract}
We record a five-parameter quadratic algebraic generating function whose
coefficient sequence has prescribed initial Hankel determinants
$(\lambda,m,r,\eta)$ and whose Hankel transform belongs to the Somos--4
family $A(1,\tau)$.  The Hankel determinants and their parameters are
introduced first, before the explicit algebraic equation.  We then describe
the two square-root branches, their coalescence at the origin, a
desingularized coefficient recurrence, the induced shift of the Hankel
parameters, the Hankel-transform identity used in the construction, and the
nonuniqueness of the coefficient sequence underlying a given Hankel transform.
\end{abstract}

\section{Hankel determinants and Somos parameters}

Let
\[
 G(z)=\sum_{n\geq0}g_nz^n
\]
be a formal power series, and for $n\geq1$ define its ordinary Hankel
determinants by
\begin{equation}
 H_n=\det(g_{i+j})_{0\leq i,j\leq n-1},
 \qquad H_0=1.
 \label{eq:Hankel}
\end{equation}

The five parameters $\lambda,m,r,\eta,\tau$ have the following roles:
\begin{equation}
 H_1=\lambda,\qquad H_2=m,\qquad H_3=r,\qquad H_4=\eta,
 \label{eq:initial-Hankel}
\end{equation}
and the Hankel sequence satisfies the Somos--4 recurrence
\begin{equation}
 H_{n+4}H_n=H_{n+3}H_{n+1}+\tau H_{n+2}^2.
 \label{eq:Somos}
\end{equation}
Thus $\tau$ is the second coefficient of the Somos family $A(1,\tau)$,
whereas $\lambda,m,r,$ and $\eta$ select the initial Hankel determinants.
The algebraic generating function producing these determinants is given in
the next section.

\section{The algebraic generating function}

Define the polynomials $A(z),B(z),C(z)$ by
\begin{align}
 A(z)={}&\lambda\eta mr^2(\tau-\eta)z^3 \\
 &+r\Bigl(-(\lambda^2\eta m+m^2(m+r^2))\tau
 +(2m-1)\lambda^2\eta^2+\eta m^2(m+r^2)
 -\lambda r(m+r^2)\Bigr)z^2 \\
 &+\lambda\Bigl(m^2r^2\tau-(m-1)\lambda^2\eta^2
 -m(m^2+2mr^2-m-r^2)\eta+\lambda r^3\Bigr)z \\
 &+\lambda^2\eta m(m-1)r,
 \label{eq:A}\\[0.4em]
 B(z)={}&-\lambda^2\eta mr^2\tau z^3 \\
 &+r\Bigl((2\lambda^3\eta m+\lambda m^2(2m+r^2))\tau
 -(2m-1)\lambda^3\eta^2-\lambda\eta m^3
 +\lambda^2r(2m+r^2)\Bigr)z^2 \\
 &+\lambda^2\Bigl(-2m^2r^2\tau+2(m-1)\lambda^2\eta^2
 +m(2m^2+2mr^2-2m-r^2)\eta-2\lambda r^3\Bigr)z \\
 &-2\lambda^3\eta m(m-1)r,
 \label{eq:B}\\[0.4em]
 C(z)={}&-\lambda^2mr\Bigl(\tau(\lambda^2\eta+m^2)+\lambda r\Bigr)z^2 \\
 &+\Bigl(\lambda^3m^2r^2\tau+\lambda^4r^3
 +\eta(m-1)\bigl(\eta(-15\lambda^4+85\lambda^3-225\lambda^2
 +274\lambda-120) \\
 &\hspace{7.5em}-\lambda^3m^2
 -3(\lambda-1)(\lambda-2)(\lambda-3)(\lambda-4)(\lambda-5)\bigr)\Bigr)z \\
 &+\lambda^4\eta m(m-1)r.
 \label{eq:C}
\end{align}

The generating function is the formal-power-series branch of
\begin{equation}
 A(z)G(z)^2+B(z)G(z)+C(z)=0
 \label{eq:quadratic}
\end{equation}
selected by the initial conditions implicit in
\eqref{eq:initial-Hankel}, beginning with $g_0=\lambda$ and $g_1=m$.

\begin{remark}[The apparently mysterious integer coefficients]
The polynomial occurring in \eqref{eq:C} is simply
\begin{align}
 -15\lambda^4+85\lambda^3-225\lambda^2+274\lambda-120
 &=\frac{6!}{\lambda}\binom{\lambda}{6}-\lambda^5 \\
 &=(\lambda-1)(\lambda-2)(\lambda-3)(\lambda-4)(\lambda-5)-\lambda^5.
 \label{eq:binomial-identity}
\end{align}
Thus the coefficients $-15,85,-225,274,-120$ arise by removing the leading term $\lambda^5$ from
the falling factorial $(\lambda-1)\cdots(\lambda-5)$.
\end{remark}

\section{The two algebraic branches}

With
\[
 D(z)=B(z)^2-4A(z)C(z),
\]
the two algebraic branches may be represented by
\begin{equation}
 G(z)=\frac{-B(z)\mathbin{\pm}\sqrt{D(z)}}{2A(z)}.
 \label{eq:explicit}
\end{equation}
When the constant term of $A$ vanishes, the rationalized expressions
\[
 \frac{-2C(z)}{B(z)\mp\sqrt{D(z)}}
\]
can be used instead.  The sign in \eqref{eq:explicit} cannot be fixed
uniformly over the full parameter space.

Write
\[
 P(z,X)=A(z)X^2+B(z)X+C(z).
\]
The constant coefficients of \eqref{eq:A}--\eqref{eq:C} are
\[
 A_0=\lambda^2\eta m(m-1)r,\qquad
 B_0=-2\lambda^3\eta m(m-1)r,\qquad
 C_0=\lambda^4\eta m(m-1)r.
\]
Consequently,
\begin{equation}
 P(0,X)=A_0(X-\lambda)^2.
 \label{eq:double-root}
\end{equation}
Both algebraic branches therefore coalesce at $X=\lambda$ when $z=0$, and
\begin{equation}
 \partial_XP(0,\lambda)=2A_0\lambda+B_0=0.
 \label{eq:vanishing-derivative}
\end{equation}
The simple-root form of the formal implicit-function theorem is not
applicable at the origin.  The desired branch is instead selected by
expanding the square root and requiring the prescribed coefficients,
beginning with $g_0=\lambda$ and $g_1=m$.

Set
\begin{equation}
 G(z)=\lambda+zF(z),\qquad
 F(z)=\sum_{q\geq0}f_qz^q,\qquad f_q=g_{q+1}.
 \label{eq:desingularization}
\end{equation}
For
\[
 D_j=2\lambda A_j+B_j,\qquad
 E_j=\lambda^2A_j+\lambda B_j+C_j,
\]
where $C_3=0$, substitution into $P(z,G(z))=0$ gives
\[
 P(z,\lambda+zF)=\sum_{j=0}^3E_jz^j
 +z\left(\sum_{j=0}^3D_jz^j\right)F
 +z^2\left(\sum_{j=0}^3A_jz^j\right)F^2.
\]
Equation \eqref{eq:double-root} gives $E_0=D_0=0$.  An ordinary power-series
branch additionally requires
\begin{equation}
 E_1=0.
 \label{eq:compatibility}
\end{equation}
The next coefficient equation is
\begin{equation}
 A_0f_0^2+D_1f_0+E_2=0.
 \label{eq:first-coefficient}
\end{equation}
Since $f_0=g_1=m$, the prescribed value of $m$ selects one of the two
branches of the desingularized quadratic.

\section{Hankel-level shift of the parameters}

To display the shift independently of the letters used above, let
\[
 G_{0;\ell,m,r,n,t}(x)=\sum_{q\geq0}g_qx^q
\]
denote the selected formal-power-series solution of \eqref{eq:quadratic}
after the substitutions
\[
 \lambda=\ell,\qquad \eta=n,\qquad \tau=t.
\]
Its Hankel determinants satisfy
\begin{equation}
 (H_1,H_2,H_3,H_4)=(\ell,m,r,n),\qquad
 H_{q+4}H_q=H_{q+3}H_{q+1}+tH_{q+2}^2.
 \label{eq:parameter-meaning}
\end{equation}
The first Somos step gives
\begin{equation}
 H_5=\frac{H_2H_4+tH_3^2}{H_1}
     =\frac{mn+tr^2}{\ell}.
 \label{eq:fifth-Hankel}
\end{equation}
Hence shifting the Hankel sequence induces the parameter map
\begin{equation}
 (\ell,m,r,n,t)\longmapsto
 \left(m,r,n,\frac{mn+tr^2}{\ell},t\right).
 \label{eq:parameter-shift}
\end{equation}

The point of this shift in the present construction is the following
Hankel-transform identity.  Remove the first coefficient of
$G_{0;\ell,m,r,n,t}$ and shift the remaining series to the origin:
\begin{equation}
 \widehat G(x)
 =\frac{G_{0;\ell,m,r,n,t}(x)-g_0}{x}
 =\sum_{q\geq0}g_{q+1}x^q.
 \label{eq:coefficient-tail}
\end{equation}
Also form the generating function with shifted parameters,
\begin{equation}
 \widetilde G(x)
 =G_{0;m,r,n,(mn+tr^2)/\ell,t}(x)
 =\sum_{q\geq0}\widetilde g_qx^q.
 \label{eq:shifted-parameter-gf}
\end{equation}
Although the two coefficient sequences are generally different, their
Hankel transforms are identical:
\begin{equation}
 \det(g_{i+j+1})_{0\leq i,j\leq q-1}
 =\det(\widetilde g_{i+j})_{0\leq i,j\leq q-1}
 \qquad(q\geq1).
 \label{eq:tail-shift-hankel-identity}
\end{equation}
Both sides therefore begin with
\[
 m,\quad r,\quad n,\quad \frac{mn+tr^2}{\ell},\ldots
\]
and continue according to the same Somos--4 recurrence with parameter $t$.
This Hankel-level identity, rather than equality of the generating functions,
is the shift property used in the design of the algebraic generating function.
In general,
\[
 \widehat G(x)\ne\widetilde G(x),
\]
which is consistent with the fact that different coefficient sequences can
have the same Hankel transform.

\section{A recurrence for the coefficients}

Using the notation of \eqref{eq:desingularization}, division of the expanded
quadratic equation by $z^2$ yields
\begin{equation}
 A(z)F(z)^2+(D_1+D_2z+D_3z^2)F(z)+E_2+E_3z=0.
 \label{eq:desingularized-quadratic}
\end{equation}
The branch condition is
\begin{equation}
 A_0m^2+D_1m+E_2=0,
 \label{eq:branch-selection}
\end{equation}
and the branch is simple precisely when
\begin{equation}
 2A_0m+D_1\ne0.
 \label{eq:nondegeneracy}
\end{equation}
Equivalently,
\begin{equation}
 (2A_0m+D_1)^2=D_1^2-4A_0E_2.
 \label{eq:branch-discriminant}
\end{equation}

Define
\begin{equation}
 T_q=\sum_{k=0}^{q}f_kf_{q-k},\qquad T_q=0\quad(q<0).
 \label{eq:Tq}
\end{equation}
Coefficient extraction from \eqref{eq:desingularized-quadratic} gives
\begin{equation}
 \sum_{j=0}^{3}A_jT_{q-j}
 +D_1f_q+D_2f_{q-1}+D_3f_{q-2}+E_{q+2}=0,
 \label{eq:convolution}
\end{equation}
where $f_q=0$ for $q<0$ and $E_j=0$ for $j>3$.  For $q\geq1$,
\begin{equation}
 T_q=2mf_q+\sum_{k=1}^{q-1}f_kf_{q-k}.
 \label{eq:Tq-split}
\end{equation}
Under condition \eqref{eq:nondegeneracy}, the coefficients are therefore generated
recursively by
\begin{equation}
 f_q=-\frac{
 A_0\displaystyle\sum_{k=1}^{q-1}f_kf_{q-k}
 +\displaystyle\sum_{j=1}^{3}A_jT_{q-j}
 +D_2f_{q-1}+D_3f_{q-2}+E_{q+2}}
 {2A_0m+D_1}.
 \label{eq:coefficient-recurrence}
\end{equation}
This quadratic convolution recurrence follows directly from the algebraic
generating-function equation.  On the exceptional locus
$2A_0m+D_1=0$, the two desingularized branches coalesce at the origin and
the next coefficient is not determined by the first-order coefficient equation;
higher-order compatibility or a Puiseux expansion may then be required.

Because $F$ is algebraic of degree at most two, it is D-finite.  Consequently,
the same coefficient sequence also satisfies a fixed-width linear recurrence with
polynomial coefficients in the index.

\section{Nonuniqueness of the coefficient sequence}

The Hankel transform does not in general determine the coefficient sequence.
The five-parameter algebraic generating function should therefore be viewed as
a representative of a larger equivalence class, rather than as the unique
generating function of a Somos orbit.  For example, the binomial transform
\begin{equation}
 g'_q=\sum_{k=0}^{q}\binom{q}{k}g_k
 \label{eq:binomial-transform}
\end{equation}
has generating function
\[
 G'(z)=\frac{1}{1-z}G\!\left(\frac{z}{1-z}\right)
\]
and preserves the Hankel determinants.  Scalar and exponential rescalings of the coefficients produce explicit
normalization factors in the determinants, and substitutions in the generating
function give further related coefficient sequences.  Equality of Hankel transforms must therefore not be replaced
by equality of generating functions without an independent argument.

No classification of these equivalent coefficient sequences is attempted
here.  The inverse problem of describing all coefficient sequences realizing
a given Hankel sequence remains open.

\end{document}
